% HELP.DBS SET COLOR 1390-1398; set-commands.md SET COLOR; PM wave-1 recovery.
\begin{XbaseEntry}{COLOR}{set}
Chooses the colors or display attributes used for text on the screen.
\Syntax
\SyntaxLine{\kw{SET COLOR TO} \opt{\xp{standard}}\opt{\kw{,}\xp{enhanced}}\opt{\kw{,}\xp{border}}\opt{\kw{,}\xp{background}}}
\Defaults
White on black for standard text, black on white for enhanced text, and black
for the border and background.
\Usage
An attribute is a foreground color followed by a slash and a background color.
The letters are N (black), B (blue), G (green), BG (cyan), R (red), RB
(magenta), GR or RG (brown), and W (white). Add + for an intense foreground
or * for blinking. U, I, and X give monochrome underline, inverse, and blank
attributes respectively.

The four positions select standard text, enhanced text, border, and background.
An omitted position keeps its previous attribute. The standard attribute is
applied to subsequent screen output; the other attributes are retained in the
session. SET COLOR TO with nothing after TO restores the defaults. A malformed
attribute raises \xmsg{10}{Syntax error.} and changes none of the attributes.
Choosing white on black does not change the result of \fn{ISCOLOR()}.
\Examples
\DotExample{examples/c-setcolor}
\SeeAlso{ISCOLOR(), SET CONSOLE}
\end{XbaseEntry}

